% PDF (xelatex): her sayfanın üstüne bölüm başlığını, altına ortalanmış sayfa
% numarasını basar. Harici paket (fancyhdr) gerektirmez; `book` sınıfının kendi
% düşük seviyeli sayfa-stili mekanizmasını kullanır, böylece derleme taşınabilir
% kalır. Kapak sayfası kapak.tex içinde \thispagestyle{empty} ile bunun dışında
% tutulur; bölüm açılış sayfaları `book` sınıfının `plain` stilini kullanır ve
% zaten alt-orta sayfa numarasını taşır.
\makeatletter

% Üstbilgide yalnızca bölüm başlığı görünsün: "Bölüm N." öneki ya da İngilizce
% "Chapter" sözcüğü eklenmesin. \sectionmark boşaltılır ki sağ mark boş kalsın.
\renewcommand{\chaptermark}[1]{\markboth{#1}{}}
\renewcommand{\sectionmark}[1]{}

% Numarasız bölümler (\chapter*: önsöz, Ek A, Ek B) \chaptermark'ı tetiklemez;
% bu yüzden üstbilgi önceki bölümün başlığında takılı kalır. \@schapter'ı
% yamayarak numarasız bölümlerin de kendi başlığını üstbilgiye yazmasını sağla.
\let\orig@schapter\@schapter
\renewcommand{\@schapter}[1]{\orig@schapter{#1}\markboth{#1}{}}

% Özel sayfa stili: üst-orta bölüm başlığı, alt-orta sayfa numarası.
\newcommand{\ps@kitap}{%
  \renewcommand{\@oddfoot}{\hfil\thepage\hfil}%
  \renewcommand{\@evenfoot}{\hfil\thepage\hfil}%
  \renewcommand{\@oddhead}{\hfil{\small\itshape\leftmark}\hfil}%
  \renewcommand{\@evenhead}{\hfil{\small\itshape\leftmark}\hfil}%
  \let\@mkboth\markboth%
}
\pagestyle{kitap}

\makeatother
